%% =============================================================
%%  TEMPLATE: Solución de Ejercicio — Economía Financiera
%%  Basado en el estilo del Ejercicio 4 (Cia Minera del Zinc)
%%
%%  CÓMO USAR:
%%  Reemplaza cada <<VARIABLE>> con el valor correspondiente.
%%  Las variables están listadas al inicio de cada bloque.
%%
%%  VARIABLES PRINCIPALES:
%%    <<CURSO>>          Ej: Econom\'ia Financiera 1
%%    <<GUIA_NUM>>       Ej: 4
%%    <<GUIA_TEMA>>      Ej: Mercado de Acciones
%%    <<EJ_NUM>>         Ej: 4
%%    <<EJ_TITULO>>      Ej: Cia Minera del Zinc
%%    <<EJ_SUBTITULO>>   Ej: Modelo de Descuento de Dividendos (DDM)
%%    <<FACULTAD>>       Ej: Facultad de Ciencias Empresariales y Econ\'omicas
%%    <<CARRERA>>        Ej: Carrera de Econom\'ia
%%    <<FECHA>>          Ej: Abril 2026
%%    <<REFLEXION_Q>>    Pregunta de la caja de reflexión
%% =============================================================

\documentclass[12pt,a4paper]{article}
\usepackage[utf8]{inputenc}
\usepackage[T1]{fontenc}
\usepackage[spanish]{babel}
\usepackage{amsmath,amssymb}
\usepackage{geometry}
\usepackage{booktabs}
\usepackage{array}
\usepackage{xcolor}
\usepackage{tcolorbox}
\usepackage{fancyhdr}
\usepackage{lmodern}
\usepackage{microtype}
\usepackage{tikz}
\usetikzlibrary{arrows.meta, positioning, decorations.pathreplacing}

\geometry{top=2.5cm, bottom=2.5cm, left=2.5cm, right=2.5cm, headheight=25pt}

% ---- Paleta de colores (no modificar) ----
\definecolor{azulul}{RGB}{0,70,127}
\definecolor{verdeclaro}{RGB}{200,230,201}
\definecolor{azulclaro}{RGB}{200,220,240}
\definecolor{rojoclaro}{RGB}{255,210,210}
\definecolor{amarilloclaro}{RGB}{255,249,196}
\definecolor{grisclaro}{RGB}{245,245,245}

\tcbuselibrary{skins, breakable}

% ---- Encabezado ----
\pagestyle{fancy}
\fancyhf{}
\lhead{\color{azulul}\textbf{<<CURSO>> -- Gu\'ia <<GUIA_NUM>>}}
\rhead{\color{azulul}Ejercicio <<EJ_NUM>>: <<EJ_TITULO>>}
\cfoot{\thepage}
\renewcommand{\headrulewidth}{1pt}
\renewcommand{\headrule}{\hbox to\headwidth{\color{azulul}\leaders\hrule height \headrulewidth\hfill}}

% ---- Cajas de colores (no modificar) ----
\newtcolorbox{cajaenunciado}{
  colback=azulclaro!40, colframe=azulul, fonttitle=\bfseries,
  title=Enunciado del Ejercicio <<EJ_NUM>>, breakable
}
\newtcolorbox{cajaintuicion}[1]{
  colback=verdeclaro!60, colframe=green!50!black, fonttitle=\bfseries,
  title={Intuici\'on: #1}, breakable
}
\newtcolorbox{cajaresultado}{
  colback=verdeclaro!80, colframe=green!40!black, fonttitle=\bfseries\large,
  title=Resultado Final, halign title=center
}
\newtcolorbox{cajaformula}{
  colback=azulclaro!20, colframe=azulul!80, fonttitle=\bfseries,
  title=F\'ormula clave, breakable
}
\newtcolorbox{cajareflexion}[1]{
  colback=amarilloclaro, colframe=orange!70!black, fonttitle=\bfseries,
  title={Reflexi\'on: #1}, breakable
}

\begin{document}

% ========================
%  PORTADA
% ========================
\begin{titlepage}
\centering
\vspace*{1cm}

{\color{azulul}\rule{\linewidth}{1.5pt}}
\vspace{0.5cm}

{\large\textbf{UNIVERSIDAD DE LIMA}\\[0.2cm]
<<FACULTAD>>\\
<<CARRERA>>}

\vspace{1.5cm}

{\LARGE\color{azulul}\textbf{<<CURSO_MAYUS>>}}\\[0.4cm]
{\Large Gu\'ia de Pr\'actica <<GUIA_NUM>> --- <<GUIA_TEMA>>}

\vspace{1.5cm}

{\color{azulul}\rule{0.6\linewidth}{0.8pt}}\\[0.5cm]
{\LARGE\textbf{Soluci\'on: Ejercicio <<EJ_NUM>>}}\\[0.3cm]
{\Large\textit{<<EJ_TITULO>>}}\\[0.3cm]
{\Large\textit{<<EJ_SUBTITULO>>}}\\[0.5cm]
{\color{azulul}\rule{0.6\linewidth}{0.8pt}}

\vspace{2cm}

%% Cuadro de datos clave — ajustar filas según el problema
\begin{tcolorbox}[colback=azulclaro!30, colframe=azulul, width=0.75\linewidth, halign=center]
\centering
\textbf{Datos clave del problema}\\[0.4cm]
\begin{tabular}{ll}
%% <<DATO_1_LABEL>>: & <<DATO_1_VALOR>> \\
%% <<DATO_2_LABEL>>: & <<DATO_2_VALOR>> \\
%% <<DATO_3_LABEL>>: & <<DATO_3_VALOR>> \\
%% Agregar/quitar filas según necesidad
\end{tabular}
\end{tcolorbox}

\vfill
{\large <<FECHA>>}

{\color{azulul}\rule{\linewidth}{1.5pt}}
\end{titlepage}

% ========================
%  ENUNCIADO
% ========================
\begin{cajaenunciado}
%%  Pegar aquí el texto del enunciado tal como aparece en la guía.
%%  Usar \textbf{} para destacar datos numéricos clave.
<<TEXTO_ENUNCIADO>>
\end{cajaenunciado}

\vspace{0.5cm}

% ========================
%  INTUICIÓN GENERAL
% ========================
\begin{cajaintuicion}{<<TITULO_INTUICION_GENERAL>>}
%%  Explicar el modelo/concepto central que se usará para resolver el ejercicio.
%%  Incluir la ecuación general si aplica.
\[
<<ECUACION_GENERAL>>
\]
<<TEXTO_INTUICION_GENERAL>>
\end{cajaintuicion}

\vspace{0.5cm}

%%  Tabla resumen de fases/etapas si aplica (opcional)
\noindent <<INTRO_TABLA_FASES>>

\begin{center}
\begin{tabular}{clcp{5.5cm}}
\toprule
\textbf{Fase} & \textbf{Per\'iodo} & \textbf{Par\'ametro} & \textbf{Raz\'on econ\'omica} \\
\midrule
%% <<FASE_1>> & <<PERIODO_1>> & <<PARAM_1>> & <<RAZON_1>> \\
%% <<FASE_2>> & <<PERIODO_2>> & <<PARAM_2>> & <<RAZON_2>> \\
%% <<FASE_3>> & <<PERIODO_3>> & <<PARAM_3>> & <<RAZON_3>> \\
\bottomrule
\end{tabular}
\end{center}

\vspace{1cm}

% ========================
%  PASO 1
% ========================
\section{Paso 1: <<TITULO_PASO_1>>}

\begin{cajaintuicion}{<<TITULO_INTUICION_PASO_1>>}
<<TEXTO_INTUICION_PASO_1>>
\end{cajaintuicion}

%%  Desarrollo del Paso 1
%%  Usar \subsection*{} para sub-partes si hay varias etapas

%% Ejemplo de subsección:
%% \subsection*{<<SUBTITULO_1A>>}
%% <<DESARROLLO_1A>>
%% \begin{align}
%%   <<ECUACION_1>> &= <<RESULTADO_1>> \\
%% \end{align}

\vspace{1cm}

% ========================
%  PASO 2
% ========================
\section{Paso 2: <<TITULO_PASO_2>>}

\begin{cajaintuicion}{<<TITULO_INTUICION_PASO_2>>}
<<TEXTO_INTUICION_PASO_2>>
\end{cajaintuicion}

\begin{cajaformula}
%%  Escribir aquí la fórmula clave de este paso
\[
<<FORMULA_CLAVE_PASO_2>>
\]
\end{cajaformula}

%%  Desarrollo del Paso 2
\begin{align}
%% <<VARIABLE_1>> &= <<EXPRESION_1>> = \mathbf{<<RESULTADO_1>>} \\[6pt]
%% <<VARIABLE_2>> &= <<EXPRESION_2>> = \mathbf{<<RESULTADO_2>>}
\end{align}

\vspace{1cm}

% ========================
%  PASO 3
% ========================
\section{Paso 3: <<TITULO_PASO_3>>}

\begin{cajaintuicion}{<<TITULO_INTUICION_PASO_3>>}
<<TEXTO_INTUICION_PASO_3>>
\end{cajaintuicion}

%%  Fórmula completa del paso 3
\[
<<FORMULA_COMPLETA_PASO_3>>
\]

%%  Tabla resumen de resultados intermedios
\begin{center}
\begin{tabular}{clrr}
\toprule
\textbf{Per\'iodo} & \textbf{Flujo} & \textbf{Valor (US\$)} & \textbf{Valor Presente (US\$)} \\
\midrule
%% Filas del cálculo...
\midrule
 & \textbf{TOTAL} & & \textbf{<<TOTAL_VP>>} \\
\bottomrule
\end{tabular}
\end{center}

\[
\boxed{<<RESULTADO_INTERMEDIO_PASO_3>>}
\]

\vspace{1cm}

% ========================
%  PASO 4  (agregar o quitar pasos según el ejercicio)
% ========================
\section{Paso 4: <<TITULO_PASO_4>>}

\begin{cajaintuicion}{<<TITULO_INTUICION_PASO_4>>}
<<TEXTO_INTUICION_PASO_4>>
\end{cajaintuicion}

\[
<<FORMULA_PASO_4>>
\]

\begin{tcolorbox}[colback=verdeclaro!80, colframe=green!40!black,
                  halign=center, boxsep=8pt]
\Large
\[
\boxed{\textbf{<<RESULTADO_FINAL_DESTACADO>>}}
\]
\end{tcolorbox}

\vspace{1cm}

% ========================
%  RESULTADO FINAL
% ========================
\begin{cajaresultado}
\begin{center}
\large
\begin{tabular}{p{7cm}p{5.5cm}}
%% <<LABEL_RESULTADO_1>>: & <<VALOR_RESULTADO_1>> \\[4pt]
%% <<LABEL_RESULTADO_2>>: & <<VALOR_RESULTADO_2>> \\[4pt]
%% <<LABEL_RESULTADO_3>>: & \textbf{<<VALOR_RESULTADO_3>>} \\
\end{tabular}
\end{center}
\end{cajaresultado}

\vspace{1cm}

% ========================
%  REFLEXIÓN
% ========================
\begin{cajareflexion}{<<REFLEXION_Q>>}
%%  Análisis adicional, caso especial, o lección clave del ejercicio.
<<TEXTO_REFLEXION>>

\textbf{Lecci\'on clave:} <<LECCION_CLAVE>>
\end{cajareflexion}

\vspace{1cm}

% ========================
%  LÍNEA DE TIEMPO  (ajustar coordenadas según número de períodos)
% ========================
\newpage
\section*{L\'inea de Tiempo del Problema}

\vspace{0.5cm}

\begin{center}
\resizebox{\linewidth}{!}{%
\begin{tikzpicture}[scale=0.95, font=\small]

%% ---- Eje temporal ----
\draw[thick, -{Stealth}] (-0.3, 0) -- (14.5, 0) node[right] {Tiempo};

%% ---- Marcas en el eje (ajustar \foreach según períodos) ----
%% Ejemplo para t=0..7+:
%% \foreach \x/\t in {0/0, 1.8/1, 3.6/2, 5.4/3, 7.2/4, 9.0/5, 10.8/6, 12.6/7+} {
%%   \draw[thick] (\x, 0.12) -- (\x, -0.12);
%%   \node[below=4pt] at (\x, 0) {$t=\t$};
%% }

%% ---- Fases (rectángulos sombreados) ----
%% \fill[green!25] (0, 0.5) rectangle (5.4, 1.6);
%% \draw[green!60!black, thick] (0, 0.5) rectangle (5.4, 1.6);
%% \node[align=center] at (2.7, 1.05) {\textbf{Fase 1}\\<<PARAM_FASE_1>>};

%% ---- Flechas de flujos ----
%% \foreach \x/\d in {1.8/<<F1>>, 3.6/<<F2>>, ...} {
%%   \draw[blue!70!black, -{Stealth}, thick] (\x, 0) -- (\x, -1.2);
%%   \node[below=2pt, blue!70!black, align=center] at (\x, -1.3) {\$\d};
%% }

%% ---- Valor terminal ----
%% \draw[purple!80!black, -{Stealth}, very thick] (<<XT>>, 0) -- (<<XT>>, -2.3);
%% \node[below=2pt, purple!80!black, align=center] at (<<XT>>, -2.4)
%%   {$P_{<<T>>}=\$<<PT>>$\\(perpetuidad)};

%% ---- Corchete "Descuentan a hoy" ----
%% \draw[thick, decorate, decoration={brace, amplitude=8pt, mirror}]
%%   (-0.1, -2.9) -- (<<XMAX>>, -2.9);
%% \node[below=10pt, align=center] at (<<XMID>>, -2.9)
%%   {Todos estos flujos se descuentan\\a $t=0$};

\end{tikzpicture}}
\end{center}

\vspace{0.8cm}

% ========================
%  VERIFICACIÓN NUMÉRICA
% ========================
\section*{Verificaci\'on num\'erica compacta}

\begin{align*}
%%  Escribir la suma compacta de VP aquí
<<VERIFICACION_SUMA>>
\end{align*}

\[
\Rightarrow\quad \boxed{\textbf{<<RESULTADO_VERIFICACION>>}}
\]

\end{document}
